\documentclass[11pt]{article}

\usepackage[margin=1in]{geometry}
\usepackage{amsmath,amssymb,amsthm}
\usepackage{enumitem}
\usepackage[hidelinks]{hyperref}

\newtheorem{theorem}{Theorem}
\newtheorem{lemma}[theorem]{Lemma}
\newtheorem{corollary}[theorem]{Corollary}
\newtheorem*{maintheorem}{Main Theorem}

\newcommand{\lca}{\operatorname{lca}}
\newcommand{\Opt}{\operatorname{Opt}}
\newcommand{\Cost}{\operatorname{Cost}}
\newcommand{\code}[1]{\texttt{#1}}

\title{The strategy computed by \code{compute\_optimal\_guesses.py} is optimal}
\author{}
\date{}

\begin{document}

\maketitle

%% ---------------------------------------------------------------------------
\section{The game}
\label{sec:game}

The game is played on a fixed finite rooted tree $T$ with root $r$. Each node $v$
has a boolean attribute $\code{canBeAnswer}(v)$. We call the nodes with
$\code{canBeAnswer}(v) = \code{true}$ \emph{answer nodes}, write $A$ for the set
of answer nodes, and let $n = |A|$, which we assume is at least~$1$. The tree is
the same in every round.

A round is played as follows.
\begin{enumerate}
  \item A hidden answer $a$ is chosen uniformly at random from $A$.
  \item The player guesses one node at a time, and every guess must be an answer
    node. After each guess $g$, the player is told whether $g = a$. If so, the
    round ends. If not, the player is also told $\lca(g, a)$, the lowest common
    ancestor of $g$ and $a$, that is, the root of the smallest subtree of $T$
    that contains both $g$ and $a$.
\end{enumerate}
The number of guesses in a round counts every guess, including the final, correct
one. The player's goal is to minimize the expected number of guesses.

Throughout, a \emph{subtree} is a node together with all of its descendants, and
every node counts as an ancestor and as a descendant of itself. So if the answer
$a$ is a descendant of the guess $g$, the player is told $\lca(g, a) = g$, and if
$g$ is a descendant of $a$, she is told $\lca(g, a) = a$.

The player may choose each guess based on everything she has been told so far in
the round, so a strategy is a decision tree: it specifies the first guess and, for
each feedback the player can get after a wrong guess, how to continue.
Section~\ref{sec:strategies} makes this precise. Since the answer is uniform on
$A$, the expected number of guesses of a strategy is $1/n$ times its
\emph{cost}: the sum, over all $a \in A$, of the number of guesses the strategy
makes when the answer is $a$. So minimizing the expected number of guesses is the
same as minimizing the cost. A strategy that makes random choices can be viewed as
first choosing a deterministic strategy at random and then following it, so its
expected number of guesses is an average of those of deterministic strategies,
and is at least the smallest of them. We therefore only consider deterministic
strategies.

%% ---------------------------------------------------------------------------
\section{The algorithm}
\label{sec:algorithm}

This section restates, in mathematical notation, what
\code{compute\_optimal\_guesses.py} computes in the functions
\code{count\_answers}, \code{sort\_options}, \code{option\_guess},
\code{strategy\_groups} and \code{build\_strategy}, and how the player uses its
output. The rest of the script reads the tree from a CSV file, sorts the groups,
and writes them to a CSV file.

\paragraph{Sizes.}
For a node $v$, let $T_v$ be the subtree rooted at $v$, let $A_v = A \cap T_v$ be
the answer nodes in it, and let $s(v) = |A_v|$. (\code{count\_answers} stores
$s(v)$ in \code{v.size}.) The \emph{height} of $v$ is the largest number of edges
on a path from $v$ down to a node of $T_v$, so a child has smaller height than its
parent.

\paragraph{Options.}
The \emph{options} of a node $v$ are the places where an answer node of $T_v$ can
be:
\begin{itemize}
  \item each child $c$ of $v$ with $s(c) \ge 1$, standing for ``the answer is in
    $T_c$'';
  \item if $v$ is an answer node, one more option, written $\hat v$, standing for
    ``the answer is $v$ itself''. (The code represents this option by $v$.)
\end{itemize}
Let $O(v)$ be the set of options of $v$, and $k_v = |O(v)|$. For an option $o$ of
$v$, let $A(o)$ be the set of answer nodes it stands for, so $A(c) = A_c$ for a
child $c$ and $A(\hat v) = \{v\}$, and let the \emph{weight} of $o$ be
$w(o) = |A(o)|$, so $w(c) = s(c)$ and $w(\hat v) = 1$ (\code{option\_size}). The
sets $A(o)$ for $o \in O(v)$ are nonempty and pairwise disjoint, and their union
is $A_v$. In particular, $v$ has at least one option if and only if $s(v) \ge 1$.

\code{sort\_options} sorts the options of $v$ into a list
\[
  L(v) = \bigl(\ell_v(0), \ell_v(1), \dots, \ell_v(k_v - 1)\bigr)
  \quad\text{with}\quad
  w(\ell_v(0)) \ge w(\ell_v(1)) \ge \dots \ge w(\ell_v(k_v - 1)).
\]
Options of equal weight are ordered alphabetically, ignoring case, by the name of
the guess that checks them (defined below). This tie-breaking rule plays no role
in the proof: everything below holds for every ordering of $L(v)$ by nonincreasing
weight. For $0 \le i \le k_v$, let
\[
  R(v, i) = \{\ell_v(i), \ell_v(i+1), \dots, \ell_v(k_v - 1)\}
\]
be the set of options from position $i$ on, so $R(v, 0) = O(v)$ and
$R(v, k_v) = \emptyset$.

\paragraph{Guesses.}
For an option $o$ of a node $v$, the guess that \emph{checks} $o$
(\code{option\_guess}) is
\[
  \gamma(o) =
  \begin{cases}
    v & \text{if } o = \hat v,\\
    \gamma(\ell_c(0)) & \text{if } o \text{ is a child } c \text{ of } v.
  \end{cases}
\]
So the guess that checks a child $c$ is the guess that checks the first option of
$c$, which \code{sort\_options} stores in \code{c.guess}. This recursive
definition is well founded, since $\ell_c(0)$ is either $\hat c$ or a child of
$c$, which has smaller height than $c$. By induction on the height,
$\gamma(o) \in A(o)$ for every option $o$: $\gamma(\hat v) = v \in A(\hat v)$, and
$\gamma(c) = \gamma(\ell_c(0)) \in A(\ell_c(0)) \subseteq A_c$.

\paragraph{Groups.}
\code{build\_strategy} returns the following groups (\code{strategy\_groups}):
\begin{itemize}
  \item a group for the root $r$, in which the player checks every option of $r$,
    in the order $\ell_r(0), \ell_r(1), \dots, \ell_r(k_r - 1)$;
  \item a group for every other node $u$ with $k_u \ge 2$, in which the player
    checks every option of $u$ but the first, in the order
    $\ell_u(1), \dots, \ell_u(k_u - 1)$.
\end{itemize}
Each group is written as one row, which lists the guesses that check its options,
in order, and gives the total weight of these options as the size of the group.

\paragraph{The strategy.}
The docstring of the script explains how the player uses the rows. Her
\emph{group} is the last lowest common ancestor she was told, or the root before
her first wrong guess, and she makes the guesses in the row of her group in
order, until she is told a different group or guesses the answer. So at each
point, the player is in a \emph{state} $(v, i)$, where $v$ is a node and
$0 \le i < k_v$: her group is $v$, and her next guess is $\gamma(\ell_v(i))$, the
guess that checks the option $\ell_v(i)$. She starts in state $(r, 0)$. If her
guess is wrong and she is told $\lambda$:
\begin{itemize}
  \item if $\lambda = v$, she moves to $(v, i + 1)$;
  \item if $\lambda \ne v$, she moves to $(\lambda, 1)$, since the row of
    $\lambda$ starts with the guess that checks $\ell_\lambda(1)$.
\end{itemize}
(The row of the root starts with the guess that checks $\ell_r(0)$ instead, but
this doesn't matter: the player's group is always an ancestor of the answer,
being the root or the lowest common ancestor of an earlier guess and the answer,
and her guesses are in the subtree of her group. So she is always told her group
or a node below it, and never the root when her group is another node.) It is
less obvious that $(v, i + 1)$ and $(\lambda, 1)$ are states, that is, that
$i + 1 < k_v$ in the first case and $k_\lambda \ge 2$ in the second (which also
means that $\lambda$ has a row); Theorem~\ref{thm:main} shows that this is always
the case.

The player's guess and her next state depend only on her current state and on the
feedback, which depends only on the guess and the answer. So for every state
$(v, i)$, even one that doesn't come up in a real game, such as $(u, 0)$ for
$u \ne r$, we can consider \emph{playing from $(v, i)$}: starting in state
$(v, i)$ and following these rules. For an answer node $a$, we say that playing
from $(v, i)$ \emph{finds} $a$ if, when the answer is $a$, the player only moves
to states and eventually guesses $a$. We then write $N_{v,i}(a)$ for the number of
guesses she makes, including the last one.

In words, the player checks the options of her group one at a time, largest
first. To check $\hat v$, she guesses $v$. To check a child $c$, she guesses
$\gamma(c)$, which is the first guess she would make playing from $(c, 0)$, that
is, if she knew that the answer is in $T_c$. If she is then told a node
$u \ne v$, she checks the options of $u$ from the second one on.

The goal of this document is to prove:

\begin{maintheorem}
  Playing from $(r, 0)$ finds the answer, whichever answer node it is, and no
  strategy has a smaller expected number of guesses.
\end{maintheorem}

This is the case $(v, i) = (r, 0)$ of Theorem~\ref{thm:main}, which shows more
generally that playing from any state $(v, i)$ is optimal for finding an answer
that is known to be in one of the options $\ell_v(i), \dots, \ell_v(k_v - 1)$.

%% ---------------------------------------------------------------------------
\section{Outline of the proof}

\begin{enumerate}
  \item A guess that can't be right gives no information, so an optimal strategy
    only makes guesses that can be right (Lemmas~\ref{lem:feedback},
    \ref{lem:useless} and~\ref{lem:bellman}).
  \item A guess in an option $o$ of a node $v$ separates the answers in $A(o)$
    from the others. For the answers outside $A(o)$, the player only learns that
    the answer is not in $A(o)$. For the answers in $A(o)$, the game goes on
    exactly as if the answer had been known to be in $A(o)$ from the start. So
    the optimal cost is the sum of the optimal costs of the options, plus a
    waiting cost: if the options are checked in the order $o_0, o_1, \dots$, each
    answer in $A(o_j)$ pays $j$ extra guesses for the checks of
    $o_0, \dots, o_{j-1}$ (Lemmas~\ref{lem:split} and~\ref{lem:order}).
  \item The waiting cost is smallest when the options are checked largest first
    (Lemma~\ref{lem:order}, by an exchange argument).
  \item The code's strategy has the same structure. In each group, it checks the
    options largest first, and after it checks a child $c$ of its group, it plays
    for the answers in $T_c$ exactly as it would play from $(c, 0)$, as if it had
    known from the start that the answer is in $T_c$. An induction then shows
    that its cost is optimal (Theorem~\ref{thm:main}).
\end{enumerate}

%% ---------------------------------------------------------------------------
\section{Candidate sets and feedback}
\label{sec:feedback}

For a node $v$ and a set $S \subseteq O(v)$ of options of $v$, let
\[
  A(v, S) = \bigcup_{o \in S} A(o)
  \qquad\text{and}\qquad
  w(S) = |A(v, S)| = \sum_{o \in S} w(o).
\]
We call a set of this form with $S \ne \emptyset$ a \emph{candidate set}. For
example, $A = A(r, O(r))$, $A_v = A(v, O(v))$ whenever $s(v) \ge 1$, and
$\{v\} = A(v, \{\hat v\})$ for every answer node $v$. For a node $v$ and
$0 \le i \le k_v$, let
\[
  X(v, i) = A(v, R(v, i)) = A(\ell_v(i)) \cup \dots \cup A(\ell_v(k_v - 1))
\]
be the set of answers in the options of $v$ from position $i$ on. It is a
candidate set if $i < k_v$, and $X(v, k_v) = \emptyset$.

For a set $X \subseteq A$, a guess $g \in A$ and a node $\lambda$, let
\[
  X^g_\lambda = \{\, a \in X \setminus \{g\} : \lca(g, a) = \lambda \,\}
\]
be the set of answers in $X$ for which the guess $g$ is wrong and the feedback is
$\lambda$. We call these sets the \emph{feedback classes} of $g$ on $X$. The
nonempty ones partition $X \setminus \{g\}$.

\begin{lemma}[Feedback]
  \label{lem:feedback}
  Let $v$ be a node.
  \begin{enumerate}[label=(\alph*)]
    \item If $o$ and $o'$ are different options of $v$, $g \in A(o)$ and
      $a \in A(o')$, then $\lca(g, a) = v$.
    \item If $c$ is a child of $v$ and $g, a \in A_c$, then $\lca(g, a) \in T_c$;
      in particular, $\lca(g, a) \ne v$.
  \end{enumerate}
  Consequently, if $X = A(v, S)$ is a candidate set, $o \in S$, and $g \in A(o)$,
  then $X^g_v = A(v, S \setminus \{o\})$, and $X^g_\lambda = A(o)^g_\lambda$ for
  every node $\lambda \ne v$.
\end{lemma}

\begin{proof}
  (a) Both $g$ and $a$ are in $T_v$, so $v$ is a common ancestor of them. If $o$
  and $o'$ are two different children $c$ and $c'$ of $v$, then $g \in T_c$ and
  $a \in T_{c'}$, and since $T_c$ and $T_{c'}$ are disjoint, no node below $v$ is
  an ancestor of both. If $o = \hat v$, then $g = v$ is an ancestor of $a$. If
  $o' = \hat v$, then $a = v$ is an ancestor of $g$. In all cases,
  $\lca(g, a) = v$.

  (b) $c$ is a common ancestor of $g$ and $a$, so $\lca(g, a)$ is a descendant of
  $c$, and $v \notin T_c$.

  For the consequence, let $a \in X \setminus \{g\}$. If $a \notin A(o)$, then
  $a \in A(o')$ for some $o' \in S \setminus \{o\}$, so $\lca(g, a) = v$ by (a).
  If $a \in A(o)$, then $o \ne \hat v$, because $A(\hat v) = \{v\}$ has only one
  element and $a \ne g$; so $o$ is a child of $v$, and $\lca(g, a) \ne v$ by (b).
  So the answers in $X \setminus \{g\}$ with feedback $v$ are exactly those in
  $X \setminus A(o) = A(v, S \setminus \{o\})$, and for $\lambda \ne v$, the
  answers with feedback $\lambda$ are exactly the answers
  $a \in A(o) \setminus \{g\}$ with $\lca(g, a) = \lambda$, which form
  $A(o)^g_\lambda$.
\end{proof}

\begin{lemma}[A guess that can't be right gives no information]
  \label{lem:useless}
  Let $X = A(v, S)$ be a candidate set and $g \in A \setminus X$. Then
  $\lca(g, a)$ is the same node for all $a \in X$. In other words, one feedback
  class of $g$ on $X$ is $X$ itself, and all others are empty.
\end{lemma}

\begin{proof}
  If $g \in T_v$, then $g \in A_v$, so $g \in A(o)$ for some option
  $o \in O(v)$, and $o \notin S$ because $g \notin X$. Every $a \in X$ lies in
  $A(o')$ for some $o' \in S$, so $o' \ne o$, and $\lca(g, a) = v$ by
  Lemma~\ref{lem:feedback}(a).

  If $g \notin T_v$, let $x = \lca(g, v)$, and let $a \in X$. Since $x$ is an
  ancestor of $v$, and $v$ is an ancestor of $a$, $x$ is a common ancestor of $g$
  and $a$. Conversely, let $y$ be a common ancestor of $g$ and $a$. Both $y$ and
  $v$ are ancestors of $a$, so one of them is an ancestor of the other. If $y$
  were a descendant of $v$, then $g \in T_y \subseteq T_v$, which is false. So
  $y$ is an ancestor of $v$, hence a common ancestor of $g$ and $v$, hence an
  ancestor of $x$. Therefore $\lca(g, a) = x$ for every $a \in X$.
\end{proof}

%% ---------------------------------------------------------------------------
\section{Strategies and the optimal cost}
\label{sec:strategies}

Let $X \subseteq A$ be nonempty; think of it as the set of answers the player
still considers possible. A \emph{strategy for $X$} is a pair
$\sigma = (g, (\sigma_\lambda)_\lambda)$ consisting of a guess $g \in A$ and, for
each node $\lambda$ with $X^g_\lambda \ne \emptyset$, a strategy $\sigma_\lambda$
for $X^g_\lambda$. This is an inductive definition, so every strategy is a finite
tree. When the answer is $a \in X$, the strategy $\sigma$ guesses $g$; if
$a \ne g$, it gets the feedback $\lambda = \lca(g, a)$, for which
$a \in X^g_\lambda$, and continues with $\sigma_\lambda$. Let $N_\sigma(a)$ be the
number of guesses $\sigma$ makes when the answer is $a$, up to and including the
guess $a$. The \emph{cost} of $\sigma$ is
\[
  \Cost(\sigma) = \sum_{a \in X} N_\sigma(a)
  = |X| + \sum_\lambda \Cost(\sigma_\lambda),
\]
where the sum is over the nodes $\lambda$ with $X^g_\lambda \ne \emptyset$: every
answer in $X$ pays for the first guess, and each $a \in X^g_\lambda$ then pays
$N_{\sigma_\lambda}(a)$ for the remaining guesses. Let
\[
  \Opt(X) = \min \{\Cost(\sigma) : \sigma \text{ is a strategy for } X\},
  \qquad
  \Opt(\emptyset) = 0.
\]
The minimum exists because costs are positive integers and every nonempty $X$
has a strategy: by induction on $|X|$, guess any $g \in X$ and continue with any
strategy for each nonempty $X^g_\lambda$, which is a subset of
$X \setminus \{g\}$.

Ways of playing as in Section~\ref{sec:game}, in which each guess may depend on
the feedback so far, give strategies in this sense. Let $X \subseteq A$ be
nonempty, and take a way of playing that finds every answer $a \in X$, with
$N(a)$ guesses. Keeping only what it does when the answer is in $X$ gives a
strategy $\sigma$ for $X$ with $N_\sigma(a) = N(a)$ for every $a \in X$: the
first guess $g$ of $\sigma$ is the first guess of the way of playing, and for each
node $\lambda$ with $X^g_\lambda \ne \emptyset$, $\sigma_\lambda$ is obtained in
the same way from what the way of playing does after the feedback $\lambda$, for
the set $X^g_\lambda$ (this is an induction on $\max_{a \in X} N(a)$). In
particular, $\sum_{a \in X} N(a) \ge \Opt(X)$. A way of playing that doesn't
always find the answer has an infinite expected number of guesses. So the smallest
possible expected number of guesses is $\Opt(A)/n$, and a way of playing that
finds every answer, with a total of $\Opt(A)$ guesses over all answers, minimizes
the expected number of guesses.

\begin{lemma}[Bellman equation]
  \label{lem:bellman}
  For every candidate set $X$,
  \[
    \Opt(X) = |X| + \min_{g \in X} \sum_\lambda \Opt(X^g_\lambda),
  \]
  where the sum is over all nodes $\lambda$ (only finitely many terms are
  nonzero, as $\Opt(\emptyset) = 0$). In particular, $\Opt(\{a\}) = 1$ for every
  answer node $a$.
\end{lemma}

\begin{proof}
  ($\le$) For every $g \in X$, the strategy that guesses $g$ and then follows a
  strategy of cost $\Opt(X^g_\lambda)$ for each nonempty $X^g_\lambda$ is a
  strategy for $X$ with cost $|X| + \sum_\lambda \Opt(X^g_\lambda)$.

  ($\ge$) Let $\sigma = (g, (\sigma_\lambda)_\lambda)$ be a strategy for $X$ with
  $\Cost(\sigma) = \Opt(X)$. If $g \notin X$, then by Lemma~\ref{lem:useless},
  $X$ itself is the only nonempty feedback class of $g$ on $X$, so
  $\Opt(X) = |X| + \Cost(\sigma_\lambda) \ge |X| + \Opt(X) > \Opt(X)$, which is
  impossible. So $g \in X$, and
  $\Opt(X) = |X| + \sum_\lambda \Cost(\sigma_\lambda) \ge
  |X| + \sum_\lambda \Opt(X^g_\lambda)$.

  For $X = \{a\}$, which is the candidate set $A(a, \{\hat a\})$, the only
  choice is $g = a$, which has no nonempty feedback classes, so
  $\Opt(\{a\}) = 1$.
\end{proof}

%% ---------------------------------------------------------------------------
\section{Splitting the cost by options}
\label{sec:split}

For an option $o$, let $C(o) = \Opt(A(o))$ be the optimal cost when the answer is
known to be in $A(o)$. So $C(\hat v) = \Opt(\{v\}) = 1$, and $C(c) = \Opt(A_c)$
for a child $c$. For a node $v$ and a set $S \subseteq O(v)$, let
\[
  F_v(S) = \Opt(A(v, S)).
\]
So $F_v(\emptyset) = 0$, $F_v(O(v)) = \Opt(A_v)$, and
$F_v(R(v, i)) = \Opt(X(v, i))$.

\begin{lemma}
  \label{lem:split}
  Let $v$ be a node, $S \subseteq O(v)$ nonempty, $X = A(v, S)$, and $o \in S$.
  Then
  \[
    \min_{g \in A(o)} \Bigl( |X| + \sum_\lambda \Opt(X^g_\lambda) \Bigr)
    = C(o) + w(S) - w(o) + F_v(S \setminus \{o\}).
  \]
  Consequently,
  \[
    F_v(S) = \min_{o \in S}
      \bigl( C(o) + w(S) - w(o) + F_v(S \setminus \{o\}) \bigr).
  \]
\end{lemma}

In words: to find an answer in $X$, the player first chooses an option $o$ to
check. The answers outside $A(o)$ pay for one guess and then continue as if $o$
had never been an option; the answers inside $A(o)$ pay exactly what they would
pay if the answer had been known to be in $A(o)$ from the start.

\begin{proof}
  Let $g \in A(o)$. By Lemma~\ref{lem:feedback}, $X^g_v = A(v, S \setminus \{o\})$
  and $X^g_\lambda = A(o)^g_\lambda$ for every node $\lambda \ne v$, so, since
  $|X| = w(S)$,
  \[
    |X| + \sum_\lambda \Opt(X^g_\lambda)
    = \bigl( w(S) - w(o) \bigr) + F_v(S \setminus \{o\})
      + \Bigl( w(o) + \sum_{\lambda \ne v} \Opt\bigl(A(o)^g_\lambda\bigr) \Bigr).
  \]
  Only the last term depends on $g$.

  If $o = \hat v$, then $g = v$ and $A(o) \setminus \{g\} = \emptyset$, so all the
  classes $A(o)^g_\lambda$ are empty, and the last term is $w(\hat v) = 1 =
  C(\hat v)$.

  If $o$ is a child $c$, then $(A_c)^g_v = \emptyset$ by
  Lemma~\ref{lem:feedback}(b), so the last term is
  $|A_c| + \sum_\lambda \Opt\bigl((A_c)^g_\lambda\bigr)$, with the sum over all
  nodes $\lambda$. By Lemma~\ref{lem:bellman} for the candidate set
  $A_c = A(c, O(c))$, the minimum of this over $g \in A_c = A(o)$ is
  $\Opt(A_c) = C(c)$.

  This proves the first claim. The second follows from Lemma~\ref{lem:bellman},
  since $X$ is the union of the sets $A(o)$ for $o \in S$, so minimizing over
  $g \in X$ is the same as minimizing over $o \in S$ and then over $g \in A(o)$.
\end{proof}

%% ---------------------------------------------------------------------------
\section{Checking the options largest first}
\label{sec:order}

\begin{lemma}
  \label{lem:order}
  Let $v$ be a node and $S \subseteq O(v)$ with $|S| = k$. Then
  \[
    F_v(S) = \sum_{o \in S} C(o)
      + \min_{(o_0, \dots, o_{k-1})} \sum_{j=0}^{k-1} j \, w(o_j),
  \]
  where the minimum is over all orderings $(o_0, \dots, o_{k-1})$ of $S$. The
  minimum is attained by every ordering with
  $w(o_0) \ge w(o_1) \ge \dots \ge w(o_{k-1})$.
\end{lemma}

\begin{proof}
  We prove the formula by induction on $k$. For $k = 0$, both sides are $0$. For
  $k \ge 1$, Lemma~\ref{lem:split} and the induction hypothesis for the sets
  $S \setminus \{o\}$ give
  \[
    F_v(S) = \min_{o \in S} \Bigl( C(o) + w(S \setminus \{o\})
      + \sum_{o' \in S \setminus \{o\}} C(o')
      + \min_{(o'_0, \dots, o'_{k-2})} \sum_{j=0}^{k-2} j \, w(o'_j) \Bigr),
  \]
  where the inner minimum is over the orderings of $S \setminus \{o\}$. Since
  $w(S \setminus \{o\}) = \sum_{j=0}^{k-2} w(o'_j)$ for each such ordering,
  \[
    F_v(S) = \sum_{o \in S} C(o) + \min_{o \in S}
      \min_{(o'_0, \dots, o'_{k-2})} \Bigl( 0 \cdot w(o)
      + \sum_{j=0}^{k-2} (j + 1) \, w(o'_j) \Bigr).
  \]
  The expression in parentheses is $\sum_{j=0}^{k-1} j \, w(o_j)$ for the
  ordering $(o_0, o_1, \dots, o_{k-1}) = (o, o'_0, \dots, o'_{k-2})$ of $S$, and
  every ordering of $S$ arises this way from exactly one choice of $o$ and
  $(o'_0, \dots, o'_{k-2})$. This proves the formula.

  Now let $(o_0, \dots, o_{k-1})$ be an ordering that minimizes
  $\sum_j j \, w(o_j)$; one exists since there are finitely many orderings. If
  $w(o_j) < w(o_{j+1})$ for some $j$, then swapping $o_j$ and $o_{j+1}$ would
  change the sum by
  \[
    \bigl( j \, w(o_{j+1}) + (j + 1) \, w(o_j) \bigr)
    - \bigl( j \, w(o_j) + (j + 1) \, w(o_{j+1}) \bigr)
    = w(o_j) - w(o_{j+1}) < 0,
  \]
  contradicting minimality. So every minimizing ordering has nonincreasing
  weights. All orderings with nonincreasing weights have the same sequence of
  weights $w(o_0), \dots, w(o_{k-1})$, and therefore the same sum, so each of
  them attains the minimum.
\end{proof}

\begin{corollary}
  \label{cor:first}
  For every node $v$ and $0 \le i < k_v$,
  \[
    F_v(R(v, i)) = C(\ell_v(i)) + w(R(v, i)) - w(\ell_v(i)) + F_v(R(v, i + 1)).
  \]
\end{corollary}

That is, when the remaining options are $R(v, i)$, it is optimal to check the
largest one, $\ell_v(i)$, next.

\begin{proof}
  The orderings $(\ell_v(i), \dots, \ell_v(k_v - 1))$ of $R(v, i)$ and
  $(\ell_v(i + 1), \dots, \ell_v(k_v - 1))$ of $R(v, i + 1)$ have nonincreasing
  weights, so by Lemma~\ref{lem:order},
  \begin{align*}
    F_v(R(v, i)) &= \sum_{j=i}^{k_v - 1} C(\ell_v(j))
      + \sum_{j=i}^{k_v - 1} (j - i) \, w(\ell_v(j)), \\
    F_v(R(v, i + 1)) &= \sum_{j=i+1}^{k_v - 1} C(\ell_v(j))
      + \sum_{j=i+1}^{k_v - 1} (j - i - 1) \, w(\ell_v(j)).
  \end{align*}
  The difference is $C(\ell_v(i)) + \sum_{j=i+1}^{k_v - 1} w(\ell_v(j))
  = C(\ell_v(i)) + w(R(v, i)) - w(\ell_v(i))$.
\end{proof}

%% ---------------------------------------------------------------------------
\section{The strategy is optimal}
\label{sec:main}

\begin{theorem}
  \label{thm:main}
  For every node $v$ and every $0 \le i < k_v$, playing from $(v, i)$ finds every
  answer $a \in X(v, i)$, and
  \[
    \sum_{a \in X(v, i)} N_{v,i}(a) = \Opt(X(v, i)).
  \]
  In particular, playing from $(r, 0)$ finds every answer in $X(r, 0) = A$, with
  a total of $\Opt(A)$ guesses, so it minimizes the expected number of guesses.
\end{theorem}

\begin{proof}
  By induction on the height of $v$, and for nodes of the same height, on
  $k_v - i$. Let $X = X(v, i)$, $S = R(v, i)$, $o = \ell_v(i)$, and
  $g = \gamma(o)$, the first guess of playing from $(v, i)$. Recall that
  $g \in A(o)$ and $\Opt(X) = F_v(S)$. The set $X$ is the disjoint union of
  $X(v, i + 1)$ and $A(o)$, and we deal with these two parts in turn.

  \emph{The answers in $X(v, i + 1)$.} Let $a \in X(v, i + 1)$, so that
  $i + 1 < k_v$ and $a \in A(o')$ for some option $o' \ne o$ of $v$. By
  Lemma~\ref{lem:feedback}(a), $\lca(g, a) = v$, so the guess $g$ is wrong, and
  the player moves to the state $(v, i + 1)$. By the induction hypothesis for
  $(v, i + 1)$, playing from $(v, i + 1)$ finds $a$. So playing from $(v, i)$
  finds $a$, with $N_{v,i}(a) = 1 + N_{v,i+1}(a)$, and by the induction
  hypothesis again,
  \[
    \sum_{a \in X(v, i + 1)} N_{v,i}(a)
    = |X(v, i + 1)| + \Opt(X(v, i + 1))
    = w(S) - w(o) + F_v(R(v, i + 1)).
  \]
  (If $i + 1 = k_v$, all three expressions are $0$.)

  \emph{The answers in $A(o)$, if $o = \hat v$.} Then $A(o) = \{v\}$ and $g = v$,
  so the only answer in $A(o)$ is found with one guess, and $1 = C(\hat v)$.

  \emph{The answers in $A(o)$, if $o$ is a child $c$ of $v$.} Since
  $s(c) \ge 1$, $c$ has an option, so $(c, 0)$ is a state, and playing from
  $(c, 0)$ also starts with the guess $\gamma(\ell_c(0)) = \gamma(c) = g$. Let
  $a \in A_c$. If $a = g$, both find $a$ with one guess. Otherwise, the feedback
  $\lambda = \lca(g, a)$ is in $T_c$ by Lemma~\ref{lem:feedback}(b), so
  $\lambda \ne v$, and playing from $(v, i)$ moves to $(\lambda, 1)$. Playing
  from $(c, 0)$ moves to $(c, 1)$ if $\lambda = c$, since $c$ is its group, and to
  $(\lambda, 1)$ otherwise; so it also moves to $(\lambda, 1)$. From then on, the
  two plays are the same, since what the player does depends only on her state
  and the answer. The height of $c$ is smaller than that of $v$, so by the
  induction hypothesis for $(c, 0)$, playing from $(c, 0)$ finds every answer in
  $X(c, 0) = A_c$ (in particular, it only moves to states, so $(\lambda, 1)$ is
  one), and $\sum_{a \in A_c} N_{c,0}(a) = \Opt(A_c) = C(c)$. Hence playing from
  $(v, i)$ finds every $a \in A_c$, with $N_{v,i}(a) = N_{c,0}(a)$, and
  $\sum_{a \in A_c} N_{v,i}(a) = C(c)$.

  In both cases, playing from $(v, i)$ finds every answer in $A(o)$, with a total
  of $C(o)$ guesses. Adding up the two parts and using
  Corollary~\ref{cor:first},
  \[
    \sum_{a \in X} N_{v,i}(a) = C(o) + w(S) - w(o) + F_v(R(v, i + 1))
    = F_v(S) = \Opt(X).
  \]
  The last claim of the theorem follows from Section~\ref{sec:strategies}, since
  playing from $(r, 0)$ is a way of playing as in Section~\ref{sec:game}.
\end{proof}

As a by-product, the proof gives a closed formula for the optimal expected number
of guesses.

\begin{corollary}[The optimal expected number of guesses]
  For a node $v$, let $\pi(v) = \sum_{j=0}^{k_v - 1} j \, w(\ell_v(j))$. Then
  \[
    \Opt(A) = n + \sum_{v \in T} \pi(v),
  \]
  so the optimal expected number of guesses is
  $1 + \frac{1}{n} \sum_{v \in T} \pi(v)$.
\end{corollary}

In words: $n$ counts the final, correct guess for each answer, and $\pi(v)$ is the
waiting cost at $v$: when the options of $v$ are checked in the order $L(v)$, the
answers in the $j$-th option (counting from $0$) wait for $j$ checks.

\begin{proof}
  We show that $\Opt(A_v) = s(v) + \sum_{u \in T_v} \pi(u)$ for every node $v$
  with $s(v) \ge 1$, by induction on the height of $v$; the case $v = r$ is the
  claim. By Lemma~\ref{lem:order} for $S = O(v)$ and the ordering $L(v)$,
  \[
    \Opt(A_v) = F_v(O(v)) = \sum_{o \in O(v)} C(o) + \pi(v).
  \]
  The options of $v$ are the children $c$ with $s(c) \ge 1$, for which
  $C(c) = \Opt(A_c) = s(c) + \sum_{u \in T_c} \pi(u)$ by the induction
  hypothesis, and, if $v \in A$, the option $\hat v$, with $C(\hat v) = 1$. The
  children $c$ with $s(c) = 0$ have $\pi(u) = 0$ for every $u \in T_c$, since
  such nodes $u$ have no options. Adding up gives
  $\Opt(A_v) = s(v) + \sum_{u \in T_v} \pi(u)$.
\end{proof}

\end{document}
